\noindent Každý, kto používa kancelársky balík, ako napríklad LibreOffice, už pravdepodobne narazil na funkciu Kontrola pravopisu a~gramatiky, ktorá poukazuje na pravopisné chyby a~navrhuje ich opravu. 40 rokov po tom, čo Ralph Gorin uviedol prvý program na kontrolu pravopisu, sa tieto programy jazyka stali oveľa sofistikovanejšími a~už nepracujú len na princípe porovnávania zoznamu vybraných slov s~pravopisným slovníkom. Oproti jazykovo závislým algoritmom na zvládnutie morfológie (napr. tvorenie plurálu) existujú aj algoritmy schopné rozpoznať syntaktické chyby, typu chýbajúce sloveso alebo sloveso nezhodné s~podmetom v~osobe a~čísle, ako to môžeme pozorovať napríklad aj vo vete ‘She *write a~letter.’ („Ona písať list.“). Najdostupnejšie funkcie kontroly pravopisu (vrátane uplatnených v~balíku LibreOffice) však v~nasledujúcej prvej strofe básne Jerrolda H. Zara založenej na homofónii nenájdu žiadnu chybu (1992)\footnote{\url{http://www.bios.niu.edu/zar/zar.shtml}}:

\begin{verse}
\emph{%
Eye have a~spelling chequer\\
It came with my Pea Sea.\\
It plane lee marks four my revue\\
Miss Steaks I~can knot sea.
}
\end{verse}

Na spracovanie tohto typu chýb je v~mnohých prípadoch potrebná analýza daného \underbar{kontextu}, ktorá je napríklad potrebná aj na rozhodnutie, či sa má isté slovo písať s~„y“ alebo s~„i“, ako napríklad v~prísloví:

\begin{verse}
\emph{%
Kto chce psa biť, palicu si nájde.\\
\smallskip
Kto chce psom byť, pána si nájde.
}
\end{verse}

Takýto postup si vyžaduje buď formuláciu gramatických
pravidiel špecifických pre daný jazyk, čo zároveň predpokladá
vysoký stupeň expertízy a~manuálnej práce, alebo využitie
takzvaného štatistického \underbar{jazykového modelu}. Takéto
modely prepočítavajú možnosť výskytu istého slova v~danom kontexte
(tzn. s~predchádzajúcimi a~nasledujúcimi slovami). Napríklad,
\emph{chce psom byť} je oveľa pravdepodobnejší sled slov ako
\emph{chce psom biť} a~naopak, \emph{chce psa biť} je oveľa
pravdepodobnejšia vetná konštrukcia než \emph{chce psa byť}
(napriek tomu by sme nepochybne dokázali vymyslieť kontexty,
v~ktorých sú gramaticky správne všetky štyri uvedené fragmenty).
Štatistický jazykový model môže byť automaticky derivovaný
využívaním veľkého množstva (korektných) jazykových dát (t.~j.
\underbar{korpusu}). Tieto prístupy však boli vyvinuté a~hodnotené
len na anglických jazykových dátach a~nedajú sa automaticky priamo
aplikovať na slovenčinu s~jej nestálym slovosledom a~bohatou flexiou.

Používanie funkcie Kontrola pravopisu a~gramatiky nie je obmedzené
len na nástroje spracovania textu, ale využíva sa aj
v~autorských systémoch. Spolu s~rastúcim počtom technických
produktov sa za posledné obdobie rapídne zvýšil aj počet technickej
dokumentácie. Strach spoločností zo sťažností zákazníkov
a~z~nárokov na náhradu škody, ktorá bola zapríčinená nesprávnymi
alebo nesprávne pochopenými inštrukciami, spôsobil, že sa
spoločnosti začali viac sústreďovať na kvalitu technickej
dokumentácie a~zároveň na medzinárodný trh. Pokroky
v~spracovávaní prirodzeného jazyka vedú k~rozvoju autorského
podporného softvéru, ktorý slúži zostavovateľovi technickej
dokumentácie na využívanie slovnej zásoby a~vetných štruktúr
v~súlade s~istými pravidlami a terminologickými
obmedzeniami.

\begin{figure*}[htb]
  \colorrule{grey3}{\textwidth}{1.5pt}
  \center
  \includegraphics[width=\textwidth]{../_media/slovak/language_checking}
  \caption{Kontrola pravopisu a gramatiky (štatistická; na báze pravidiel)}
  \label{fig:langcheckingaarch_sk}
  \colorrule{grey3}{\textwidth}{1.5pt}
\end{figure*}

\boxtext{Funkcie kontroly pravopisu a~gramatiky pre slovenský jazyk
sú väčšinou založené na slovníku základných slovných tvarov
(lem) a súbore pravidiel na odvodenie ostatných tvarov}

Existujúce zariadenia kontroly pravopisu a~gramatiky pre slovenský
jazyk sú väčšinou založené na slovníku základných slovných
tvarov (lem) skombinovanom so súborom morfologických pravidiel, ktorý
umožňuje analýzu alebo generovanie všetkých (správnych) slovných
tvarov. Hoci sa zdá tento jednoduchý uspokojivý, má dve
zásadné nevýhody. Prvou nevýhodou je nesprávne určenie zdanlivo
správnych slovných tvarov v~dôsledku nesprávneho kontextu. Druhou
nevýhodou je neschopnosť rozlišovať skutočné pravopisné chyby od
správnych slovných tvarov, ktoré však nie sú obsiahnuté
v~slovníku. Takéto slová však budú vzhľadom na prirodzené
pribúdanie nových slov, vedeckých a~technických termínov
v~lexikóne existovať stále.

Okrem kontroly pravopisu a~autorizovanej podpory je funkcia kontrola
pravopisu a~gramatiky takisto dôležitá v~oblasti výučby jazyka. Aplikácie na kontrolu gramatiky a pravopisu taktiež dokážu pri preklepoch navrhnúť správne slovo, napríklad Google frázou „Mali ste na mysli\dots“

